\begin{center}
  \section{\capitalisewords{Pi Digits in Huffman Coding: Entropy, Compression, and Randomness}}
  Adrishikhar Chowdhury, 3$^{rd}$ Year, ECE
\end{center}

\setcounter{figure}{0}
\setcounter{table}{0}

\begin{multicols}{2}
  \begin{abstract}
    \bfseries
    This report uses the first $N=10^5$ decimal digits of $\pi$ as a test bed for lossless source coding and empirical randomness. If the digits behave like independent, uniformly distributed symbols from the alphabet $\{0,1,\ldots,9\}$, their Shannon entropy is $\log_2(10)\approx3.3219$ bits per digit. A frequency analysis, chi-squared uniformity test, and Huffman-code construction can therefore compare observed statistics with this theoretical model. The report examines fixed four-bit coding, static and adaptive Huffman coding, average code length, compression ratio, implementation complexity, and round-trip decoding. A theoretical audit is also included: for ten exactly uniform symbols, an ordinary binary symbol-by-symbol Huffman code has average length $3.4$ bits per digit, not $3.3219$. Reaching the entropy limit requires arithmetic coding, range coding, or sufficiently large block coding. This distinction is essential when interpreting compression results as evidence of randomness.
  \end{abstract}

  \subsection*{\centering\uppercase{Introduction}}
  The decimal expansion of $\pi$ is widely conjectured to be normal: every digit should occur asymptotically with probability $0.1$, and every finite digit block should occur with its expected frequency. Normality has not been proved for $\pi$, but large computed prefixes show strong empirical agreement with uniform digit frequencies [3].

  \noindent
  This makes the digit stream an interesting source-coding experiment. A fixed representation assigns four bits to each decimal digit, even though ten symbols theoretically require only $\log_2(10)$ bits on average. Variable-length or arithmetic codes can reduce this redundancy, while deviations in achieved compression may reveal nonuniform frequencies, finite-sample effects, code constraints, or implementation overhead.

  \subsection*{\centering\uppercase{Objectives}}
  The analysis has four objectives:
  \begin{enumerate}\justifying
    \item estimate digit frequencies and empirical entropy;
    \item test whether the observed counts are consistent with a uniform distribution;
    \item construct and verify a lossless Huffman encoder and decoder;
    \item compare average length, storage, and runtime with theoretical and fixed-length baselines.
  \end{enumerate}

  \subsection*{\centering\uppercase{Information-Theory Model}}
  Let $p_i$ be the probability of digit $i$. The source entropy is
  \[
    H(X)=-\sum_{i=0}^{9}p_i\log_2p_i.
  \]
  For an exactly uniform decimal source, $p_i=0.1$, hence
  \[
    \begin{aligned}
      H(X)&=-10(0.1)\log_2(0.1)\\
          &=\log_2(10)\\
          &\approx3.3219\ \text{bits/digit}.
    \end{aligned}
  \]
  Entropy is a lower bound on the average length of any uniquely decodable binary code. It is not automatically the achieved length of a symbol-by-symbol Huffman code.

  \subsubsection*{Fixed-Length Baseline}
  A decimal digit can be stored in a four-bit nibble, using binary values \texttt{0000} through \texttt{1001}. The cost is
  \[
    L_{\mathrm{fixed}}=4\ \text{bits/digit},
  \]
  which is approximately $0.6781$ bits above the uniform-source entropy.

  \subsubsection*{Huffman Coding}
  Huffman coding repeatedly merges the two least frequent nodes in a minimum-priority queue. Assigning 0 and 1 to the branches creates a prefix-free tree in which more probable symbols receive shorter codewords. Its average length is
  \[
    L=\sum_{i=0}^{9}p_i\ell_i,
  \]
  where $\ell_i$ is the codeword length for digit $i$. Huffman coding guarantees
  \[
    H(X)\leq L < H(X)+1.
  \]

  \subsubsection*{Uniform Ten-Symbol Constraint}
  For ten equally likely symbols, the optimal binary Huffman tree contains six codewords of length 3 and four of length 4. The Kraft equality is
  \[
    6\left(2^{-3}\right)+4\left(2^{-4}\right)=1,
  \]
  giving
  \[
    L=\frac{6(3)+4(4)}{10}=3.4\ \text{bits/digit}.
  \]
  Therefore, a result of $3.322$ bits per digit cannot come from ordinary digit-wise binary Huffman coding unless the measurement represents entropy rather than code length, uses blocks of digits, or applies a different entropy coder.
\end{multicols}

\begin{table}[tb]
\centering
\caption{Theoretical comparison of coding approaches for nearly uniform decimal digits.}
\scriptsize
  \setlength{\tabcolsep}{3pt}
  \renewcommand{\arraystretch}{1.4}
  \begin{tabular}{@{}>{\RaggedRight\arraybackslash}p{0.25\textwidth}
                      >{\Centering\arraybackslash}p{0.17\textwidth}
                      >{\Centering\arraybackslash}p{0.18\textwidth}
                      >{\RaggedRight\arraybackslash}p{0.32\textwidth}@{}}
    \toprule
    \textbf{Representation} & \textbf{Bits/Digit} & \textbf{Relative Size} & \textbf{Interpretation} \\
    \midrule
    Fixed four-bit code & $4.000$ & $100\%$ & Simple baseline with unused binary patterns \\
    Uniform-source entropy & $3.3219$ & $83.05\%$ & Shannon lower bound, not a digit-wise prefix code \\
    Binary Huffman, uniform digits & $3.400$ & $85.0\%$ & Optimal symbol-by-symbol binary prefix code \\
    Arithmetic or range coding & Approaches $3.3219$ & Approaches $83.05\%$ & Can approach entropy for long sequences \\
    \bottomrule
  \end{tabular}
\end{table}
\vspace{-0.18em}

\begin{multicols}{2}
  \subsection*{\centering\uppercase{Uniformity Testing}}
  Table~1 compares theoretical coding approaches for nearly uniform decimal digits.
  For $N$ digits, let $f_i$ be the observed count of digit $i$. Under the uniform hypothesis, the expected count is $e_i=N/10$. The chi-squared statistic is
  \[
    \chi^2=\sum_{i=0}^{9}\frac{(f_i-e_i)^2}{e_i}.
  \]
  It is compared with a chi-squared distribution having $9$ degrees of freedom. A large $p$-value means the data do not provide sufficient evidence to reject uniformity; it does not prove independence or normality. Testing single-digit counts cannot detect every form of local correlation, so digit-pair and longer-block tests are useful extensions.

  \subsection*{\centering\uppercase{Methodology}}
  \subsubsection*{Digit Generation}
  The first $10^5$ digits after the decimal point are generated using arbitrary-precision arithmetic, such as \texttt{mpmath}, or read from a verified digit file. The leading integer part and decimal point are excluded. A checksum or independent source should be used to validate the input.

  \subsubsection*{Frequency and Entropy Analysis}
  Each digit count is measured and converted to an empirical probability
  \[
    \hat{p}_i=\frac{f_i}{N}.
  \]
  The empirical entropy is then
  \[
    \hat{H}=-\sum_{i=0}^{9}\hat{p}_i\log_2\hat{p}_i.
  \]
  Near-uniform counts should produce a value close to $\log_2(10)$.

  \subsubsection*{Static Huffman Construction}
  A minimum heap is initialised with ten leaf nodes weighted by $f_i$. The two nodes with the smallest weights are removed, merged, and returned to the heap until a single tree remains. Tree traversal assigns a binary codeword to each digit.

  \subsubsection*{Adaptive Coding}
  An adaptive Huffman encoder updates its model as digits arrive, eliminating the need to transmit a complete frequency table in advance. It can follow changing local statistics, but it has additional update complexity and still obeys integer codeword-length constraints.

  \subsubsection*{Verification and Metrics}
  The encoded stream is decoded and compared byte-for-byte with the original digit sequence. The primary metrics are
  \[
    L=\frac{B_{\mathrm{encoded}}}{N}, \qquad
    R=\frac{B_{\mathrm{encoded}}}{4N}\times100\%,
  \]
  where $B_{\mathrm{encoded}}$ is the number of coded payload bits. A complete storage comparison must also include the codebook, tree, length metadata, padding, and file headers.

  \subsection*{\centering\uppercase{Algorithm Outline}}
  \begin{enumerate}\justifying
    \item Generate or load the verified $\pi$-digit sequence.
    \item Count occurrences of digits 0 through 9.
    \item Compute empirical probabilities and entropy.
    \item Apply the chi-squared uniformity test.
    \item Build the Huffman tree from the counts.
    \item Encode every digit using the derived codebook.
    \item Decode the bitstream and verify exact recovery.
    \item Record payload length, total file size, construction time, encoding time, and decoding time.
  \end{enumerate}

  \subsection*{\centering\uppercase{Results and Critical Discussion}}
  The uploaded report states that the first $10^5$ digits are close to uniformly distributed and that chi-squared $p$-values remain above $0.05$ across several prefix lengths. This is consistent with the general empirical behaviour reported for large computed prefixes of $\pi$ [3].

  \noindent
  It also reports $3.322$ bits per digit, an $83\%$ size relative to four-bit coding, and runtimes of approximately $0.4$ seconds for Huffman coding versus $0.1$ seconds for the fixed baseline. The ratio calculation is internally consistent because $3.322/4\approx0.8305$. However, the code length is inconsistent with ordinary binary digit-wise Huffman coding for a near-uniform ten-symbol alphabet. The expected Huffman result should be near $3.4$ bits per digit, or roughly $85\%$ of the four-bit baseline.

  \noindent
  This discrepancy does not invalidate the experiment; it identifies a measurement or labelling issue. The reported $3.322$ may be empirical entropy, an arithmetic-coding result, a block-code estimate, or a payload calculation that did not measure the actual prefix-coded stream. Reproducing the experiment should print each digit's codeword length and verify that
  \[
    B_{\mathrm{encoded}}=\sum_{i=0}^{9}f_i\ell_i.
  \]

\end{multicols}

\begin{table}[b]
\centering
\caption{Recommended measurements for a reproducible compression experiment.}
\scriptsize
  \setlength{\tabcolsep}{3pt}
  \renewcommand{\arraystretch}{1.25}
  \begin{tabular}{@{}>{\bfseries\RaggedRight\arraybackslash}p{0.22\textwidth}
                      >{\RaggedRight\arraybackslash}p{0.32\textwidth}
                      >{\RaggedRight\arraybackslash}p{0.38\textwidth}@{}}
    \toprule
    \textbf{Measurement} & \textbf{Purpose} & \textbf{Reporting Requirement} \\
    \midrule
    Digit counts & Test first-order uniformity & Report all ten counts and $N$ \\
    Entropy $\hat{H}$ & Estimate ideal first-order rate & State log base and probability model \\
    Huffman payload bits & Measure actual prefix-code length & Compute $\sum f_i\ell_i$ \\
    File overhead & Measure practical storage & Include codebook, padding, and headers \\
    Round-trip test & Establish losslessness & Compare decoded and original sequences exactly \\
    Runtime & Compare computational cost & Separate generation, construction, encode, and decode time \\
    \bottomrule
  \end{tabular}
\end{table}
\vspace{-0.18em}

\begin{multicols}{2}
  \subsection*{\centering\uppercase{Implementation Considerations}}
  {\justifying
  A practical Python implementation can use \texttt{mpmath} for high-precision digit generation, \texttt{numpy.bincount} for digit frequencies, \texttt{heapq} for the Huffman priority queue, and \texttt{scipy.stats.chisquare} for the statistical test [4--6].\par}

  \begin{center}
  \begin{tcolorbox}[
    colback=black!2,
    colframe=IEEEblue!70!black,
    boxrule=0.5pt,
    arc=0mm,
    left=2mm,right=2mm,top=2mm,bottom=2mm
  ]
    \scriptsize\ttfamily\justifying
    counts = bincount(digits, minlength=10)\par
    probs = counts / len(digits)\par
    entropy = -sum(probs * log2(probs))\par
    tree = build\_huffman\_tree(counts)\par
    bits = encode(digits, tree)\par
    assert decode(bits, tree) == digits
  \end{tcolorbox}
  \end{center}

  For only ten symbols, tree construction is negligible. Encoding and decoding are $O(N)$ once the codebook is available. Digit generation and conversion may dominate runtime, so timing should separate data preparation from coding.

  \subsection*{\centering\uppercase{Randomness and Compressibility}}
  Compression is a useful diagnostic, but it is not a complete randomness test. A sequence can have uniform single-symbol frequencies and still contain predictable pairs or repeated blocks. Conversely, a finite random sequence may show small frequency imbalances. Stronger analysis should examine:
  \begin{itemize}\justifying
    \item pair and block frequencies;
    \item serial correlation and runs;
    \item entropy rate rather than single-symbol entropy;
    \item arithmetic-coding length with a contextual model;
    \item standard suites such as NIST statistical tests, with careful interpretation.
  \end{itemize}

  \subsection*{\centering\uppercase{Future Work}}
  Arithmetic coding or range coding should be used to approach the $3.3219$-bit entropy limit without integer codeword restrictions. Block Huffman coding over pairs of digits creates an alphabet of 100 symbols and reduces per-digit redundancy, although it requires a larger codebook. Context models can test whether knowledge of previous digits improves prediction; any stable improvement would indicate exploitable dependence in the tested prefix.

  \subsection*{\centering\uppercase{Conclusion}}
  Table~2 lists recommended measurements for a reproducible compression experiment.
  The decimal digits of $\pi$ provide a useful and engaging data source for studying entropy, prefix codes, statistical uniformity, and lossless compression. The first $10^5$ digits are expected to have empirical entropy close to $\log_2(10)$, and frequency tests generally support a near-uniform first-order model.

  \noindent
  The experiment also demonstrates why theoretical interpretation matters. Shannon entropy is a lower bound, whereas binary Huffman codewords have integer lengths. For ten uniform symbols, the optimal symbol-by-symbol Huffman average is $3.4$ bits per digit. Claims of $3.322$ bits per digit should therefore be identified as entropy, arithmetic coding, block coding, or re-examined as a measurement error.

  \noindent
  With transparent metrics, complete overhead accounting, and exact decode verification, the project becomes a strong demonstration of the relationship between information theory, algorithms, probability, and the numerical properties of $\pi$.

  \subsection*{\centering\uppercase{References}}
  {\small
  \begin{justify}
    \begin{enumerate}[label={\relax[\arabic*]},leftmargin=*,labelsep=0.5em,itemsep=0.20em,parsep=0pt]
      \item C. E. Shannon, “A Mathematical Theory of Communication,” \textit{Bell System Technical Journal}, vol. 27, pp. 379--423 and 623--656, 1948.
      \item D. A. Huffman, “A Method for the Construction of Minimum-Redundancy Codes,” \textit{Proceedings of the IRE}, vol. 40, no. 9, pp. 1098--1101, 1952.
      \item Y. Kanada and T. Kuroda, “Distribution of Leading Digits and Properties of the Decimal Expansion of $\pi$,” \textit{The Mathematical Intelligencer}, vol. 24, no. 1, pp. 56--60, 2002.
      \item F. Johansson et al., “mpmath: Python library for arbitrary-precision floating-point arithmetic,” Version 1.3.0.
      \item S. van der Walt, S. C. Colbert, and G. Varoquaux, “The NumPy Array: A Structure for Efficient Numerical Computation,” \textit{Computing in Science \& Engineering}, vol. 13, no. 2, pp. 22--30, 2011.
      \item P. Virtanen et al., “SciPy 1.0: Fundamental Algorithms for Scientific Computing in Python,” \textit{Nature Methods}, vol. 17, pp. 261--272, 2020.
    \end{enumerate}
  \end{justify}}

\end{multicols}

\begin{tcolorbox}[
  colback=IEEEorange!7,
  colframe=IEEEorange!75!black,
  boxrule=0.6pt,
  arc=0mm,
  left=3mm,right=3mm,top=1mm,bottom=1mm
]
\footnotesize
{\sffamily\bfseries\color{IEEEorange!75!black}DID YOU KNOW?}\quad
An optimal Huffman tree for ten equally likely symbols assigns six digits 3-bit codes and four digits 4-bit codes, averaging exactly $3.4$ bits per digit. Arithmetic or sufficiently large block coding is required to approach the $3.3219$-bit entropy limit.
\end{tcolorbox}